\section{An induced spin partner and its exact scope}
\label{sec:s5-induced-spin-scope}

\paragraph{Review status.}
The complete post-writing author reading and two independent
whole-proof readings have passed. The source interfaces have
their separately recorded scope. This is not a gate reduction.

Retain Calegari's finite central cover \(G_N\) of \(S_5\),
its index-two subgroup \(H_N\), the two-dimensional
representation \(\varrho\) of \(H_N\), and the real
quadratic field \(E\). Write \(\iota\) for the nontrivial
conjugation over \(\mathbb Q\). The known automorphic
representation \(\pi_E\) corresponds to \(\varrho\),
and the actual descended cuspidal representation \(\Pi\)
satisfies
\[
 \operatorname{BC}_{E/\mathbb Q}\Pi
       \simeq\operatorname{Sym}^3\pi_E.
 \tag{MA-RB.310}
\]
This does not identify \(\Pi\) with the finite-cover
candidate \(\xi\). The latter satisfies only the already
proved representation identity
\(\xi|_{H_N}\simeq\operatorname{Sym}^3\varrho\).
Put
\[
 \tau=\operatorname{AI}_{E/\mathbb Q}\pi_E,
 \qquad R_T=W\oplus W\eta,
 \qquad L_T=L_WL_{W\eta}.
 \tag{MA-RB.311}
\]
The notation \(R_T,L_T\) is the actual entire Hecke factor
from MA-RB-013; \(X\) is the established six-dimensional
Artin representation and its strongly corresponding
cuspidal automorphic representation is already available.

\paragraph{MA-RB-016 (the induced spin comparison is an old product).}
There is a finite-order character \(\delta\) of \(G_N\)
whose restriction to \(H_N\) is \(\det\varrho\).
Regard it also as the associated Hecke character of
\(\mathbb Q\). With the usual local Langlands normalization,
the actual Rankin--Selberg function has the exact identity
\[
 L\bigl(s,(\Pi\otimes\delta^{-2})\times\tau\bigr)
                         =L_T(s)L_X(s).
 \tag{MA-RB.312}
\]
The local identity holds at every finite place and at infinity.
Its finite conductor is \(N_TN_X\),
its archimedean factor is
\(\Gamma_{\mathbb R}(s)^8\Gamma_{\mathbb R}(s+1)^8\),
and therefore its standard completed function is exactly
\[
 \Lambda\bigl(s,(\Pi\otimes\delta^{-2})\times\tau\bigr)
                         =\Lambda_T(s)\Lambda_X(s).
 \tag{MA-RB.313}
\]
Consequently this partner produces no exceptional-prime
comparison factor and no additional real-zero holomorphy
inequality. In the residual branch of MA-RB-002 its order is
\[
 \operatorname{ord}_{\gamma}
     L\bigl(s,(\Pi\otimes\delta^{-2})\times\tau\bigr)
                       =n+k+x_1+e\ge k.
 \tag{MA-RB.314}
\]
No zero-location premise \(\mathsf A\), \(M\), or
\((\mathcal R_B)\) is used.

\begin{proof}
Calegari's central-cover lemma states that a character of
\(H_N\) is determined by its restriction to the central
cyclic group \(\Delta\). Conjugation by \(G_N\) fixes
\(\Delta\), so \(d=\det\varrho\) is invariant under
that conjugation. Choose \(t\in G_N\setminus H_N\) and
a root of unity \(z\) with \(z^2=d(t^2)\). Define
\(\delta(h)=d(h)\) for \(h\in H_N\) and
\(\delta(th)=z d(h)\). The conjugation invariance and
the relation \(t^2\in H_N\) verify multiplicativity in
the four possible coset products. This gives the asserted
finite-order character. The two extensions differ by
\(\eta\), and hence \(\delta^{-2}\) is independent of
this choice.

The two finite representations \(\varrho\) and
\(\iota\varrho\) are nonisomorphic. For example their
normalized symmetric squares are the two different
three-dimensional representations of \(A_5\), as in
Calegari's group calculation and his alternate construction
of \(X\). Thus \(\pi_E\) and \(\pi_E^\iota\) are
nonisomorphic: an isomorphism would equate their good
Frobenius characters and then their finite representations
by Chebotarev. Cyclic automorphic induction consequently
makes \(\tau\) a unitary cuspidal representation of
\(\operatorname{GL}_4/\mathbb Q\). Only its existence
and its local parameters will be needed below.

We first record the finite representation computation.
The standard two-dimensional tensor identity and the
already fixed \(A_5\) constituents give
\[
 \begin{aligned}
 (\operatorname{Sym}^3\varrho\otimes\varrho)d^{-2}
   &\simeq
       \operatorname{Sym}^4\varrho\,d^{-2}
          \oplus\operatorname{Sym}^2\varrho\,d^{-1},\\
 \operatorname{Ind}_{H_N}^{G_N}
       (\operatorname{Sym}^4\varrho\,d^{-2})
   &\simeq W\oplus W\eta,\\
 \operatorname{Ind}_{H_N}^{G_N}
       (\operatorname{Sym}^2\varrho\,d^{-1})
   &\simeq X.
 \end{aligned}
 \tag{MA-RB.315}
\]
The first identity can also be checked before the common
\(d^{-2}\) twist on \(\operatorname{diag}(a,b)\): the weights
\(a^4,b^4\) occur once and \(a^3b,a^2b^2,ab^3\)
occur twice, exactly as in the two indicated summands.
The second identity is induction of \(W|_{H_N}\);
the third is the known induced three-dimensional
\(A_5\) constituent of \(X\). In particular, projection
of tensor product through induction gives
\[
 (\xi\otimes\operatorname{Ind}_{H_N}^{G_N}\varrho)
                    \delta^{-2}\simeq R_T\oplus X.
 \tag{MA-RB.316}
\]
The central character cancels in this tensor, which is
therefore an honest \(S_5\) representation.

The decisive point is that the same calculation applies to
the actual local parameters of \(\Pi\), without giving
\(\Pi\) a global Galois representation. At a place \(v\)
where \(E_v/\mathbb Q_v\) is a field, let \(\phi_v\)
be the Weil--Deligne parameter of \(\Pi_v\). Exact local
compatibility of cyclic base change and the symmetric-cube
transfer gives
\[
 \phi_v|_{W'_{E_v}}\simeq
                         \operatorname{Sym}^3\varrho_v.
 \tag{MA-RB.317}
\]
The parameter of \(\tau_v\) is
\(\operatorname{Ind}_{W'_{E_v}}^{W'_{\mathbb Q_v}}
\varrho_v\). The projection formula for these actual
parameters, followed by (MA-RB.315), identifies
\[
 (\phi_v\otimes
    \operatorname{Ind}_{W'_{E_v}}^{W'_{\mathbb Q_v}}
        \varrho_v)\delta_v^{-2}
       \simeq (R_T\oplus X)|_{W'_{\mathbb Q_v}}.
 \tag{MA-RB.318}
\]
Indeed, restriction to the subgroup is all that the
projection formula uses. At a split place the same formula
is the direct sum over the two places above \(v\), so it
holds there as well. This also covers infinity, since
\(E\) is real quadratic and splits at the real place of
\(\mathbb Q\). In particular no ramified factor is
deleted and no unknown descent sign is selected.

The local Langlands correspondence identifies the local
Rankin--Selberg factors and epsilon factors with those of
this tensor parameter. Thus (MA-RB.318) proves (MA-RB.312)
and additivity of its conductor. At the specified double
transposition, \(R_T\) has plus/minus dimensions \((6,4)\)
and \(X\) has dimensions \((2,4)\). Their sum proves
the asserted gamma factor and (MA-RB.313).

Both \(L_T\) and \(L_X\) already have entire completed
functions. Their factors are units relative to completion
at \(\gamma\in(1/2,1)\). The checked residual orders are
\(\operatorname{ord}_{\gamma}L_T=n+k+x_1\) and
\(\operatorname{ord}_{\gamma}L_X=e\), with
\(n,x_1,e\ge0\). This proves (MA-RB.314). The formal
genuine-character assignment \(k(V\eta+W)\) gives
order exactly \(k\), so this holomorphy test does not
contradict it. This last observation is a feasibility
check, not an arithmetic zero example.
\end{proof}

\paragraph{Exact source interfaces and scope.}
The finite cover, the constituents in (MA-RB.315), and
the actual descent (MA-RB.310) are Calegari's \S\S2--4;
the supplied \texttt{S5ArtinWeb.tex}, lines194--397 and
434--556, gives exact locators. The Sasaki automorphy
input is the same previously audited input; it is not
reproved here.

Henniart, G. (2012). Induction automorphe globale pour
les corps de nombres. \emph{Bulletin de la Soci\'et\'e
Math\'ematique de France}, 140(1), 1--17,
\url{https://doi.org/10.24033/bsmf.2622}.
Theorems~2--5, pp.~5--7,
and \S\S1.11--1.12 supply cyclic induction, descent,
and the Weil--Deligne conventions. The terminal proofs
in \S2.6 and \S\S3.1--3.7, pp.~11--16, were read;
the final p.~16 remark explicitly includes infinity.
This is not a fresh audit of the imported Arthur--Clozel
trace formula or the entire local Langlands theorem.

Kim, H.~H., \& Shahidi, F. (2002). Functorial products
for \(\operatorname{GL}_2\times\operatorname{GL}_3\)
and the symmetric cube for \(\operatorname{GL}_2\).
\emph{Annals of Mathematics}, 155(3), 837--893,
\url{https://doi.org/10.2307/3062134}.
Theorem~B and its definition, pp.~838--839,
define the symmetric cube at every place. Theorem~6.1,
Proposition~6.3, Corollary~6.4, and Remark~6.6,
pp.~870--873, provide the exact all-place form used here.
The full terminal nonmonomial argument \S6.2 through
Remark~6.6 was read, including its local-component
cancellation; the preceding converse-theorem machinery
is retained as a classical input.

The finite objects in (MA-RB.315)--(MA-RB.316) are
representations of an actual finite cover in the already
fixed SGA~1, Expos\'e~V, \S7, and Deligne (1979),
Definition~6.1 and \S\S6.2--6.3 conventions. Equation
(MA-RB.318) is a local parameter identity, not a construction
of the missing global representation of \(\Pi\).
The current arithmetic gate remains MA-RB-014. No general
impossibility of other spin-cover comparisons is asserted.
